% Generated draft: requires mathematical and source review.
\documentclass[11pt]{article}
\usepackage[margin=25mm]{geometry}
\usepackage{fontspec}
\setmainfont{Latin Modern Roman}
\usepackage{amsmath,amssymb,mathtools,unicode-math}
\setmathfont{Latin Modern Math}
\usepackage{enumitem,needspace,xurl,hyperref}
\setlength{\parindent}{0pt}
\setlength{\parskip}{5pt}
\setlength{\emergencystretch}{3em}
\newcommand{\N}{\mathbb{N}}
\newcommand{\Z}{\mathbb{Z}}
\newcommand{\Q}{\mathbb{Q}}
\newcommand{\R}{\mathbb{R}}
\newcommand{\C}{\mathbb{C}}
\newcommand{\F}{\mathbb{F}}
\newcommand{\A}{\mathbb{A}}
\newcommand{\PP}{\mathbb{P}}
\newcommand{\E}{\mathbb{E}}
\newcommand{\eps}{\varepsilon}
\newcommand{\dd}{\,\mathrm{d}}
\begin{document}
\section*{UnsolvedMath ID: 11000088}
\textbf{Source problem code:} catalogue entry
\subsection{Definitions and mathematical statement}
In the Torelli-space setting of this chapter, $\mathcal T_g$ denotes the Torelli space of homologically marked genus-$g$ curves, not the contractible fully marked Teichm\"uller space. In the usual closed-surface model it is the quotient of Teichm\"uller space by the Torelli group. Let $\sigma$ be the involution induced by negating the homology marking, corresponding to $-I$ in the symplectic group. The exact space/marking conventions should be confirmed against the chapter, since they are not restated in this excerpt.

Set $R=\Z[1/2]$. Since $2$ is invertible, the involution splits homology into eigenspaces
\[
 H_\bullet(\mathcal T_g;R)=H_\bullet(\mathcal T_g;R)^+\oplus H_\bullet(\mathcal T_g;R)^-,\qquad
 H_\bullet(\mathcal T_g;R)^-=\{v:\sigma_*v=-v\},
\]
using the projector $(1-\sigma_*)/2$. This convention is supported by the adjacent source item's explicit ``invariant homology'' discussion. The extracted bullet typography does not consistently preserve upper versus lower indices, so the homology reading must still be verified rather than silently exchanged with cohomology.

\textbf{Problem 4.2.} Is $H_\bullet(\mathcal T_g;\Z[1/2])^-$ always a finitely generated $\Z[1/2]$-module? The natural reading is finite generation of the total graded module
\[
 \bigoplus_{j\ge0}H_j(\mathcal T_g;\Z[1/2])^-,
\]
which is stronger than requiring finite generation separately in every degree. The genus range and intended graded-module convention are not explicitly stated. The accompanying qualitative question asks whether the infinite topology of $\mathcal T_g$ comes from the Jacobian locus $\mathcal J_g$, rather than from the anti-invariant part.

The source suggests understanding the compact-type objects written $\mathcal T_g^c$ or $\mathcal S_g^c$. It states that $\mathcal J_g^c$ is Stein and hence $\mathcal S_g^c$ is Stein, then invokes Hamm [9] (also [5, p. 152]) for a CW model of $\mathcal S_g^c$ of dimension at most $3g-3$. It warns that $\mathcal S_g^c$ is not a rational homology manifold for $g\ge3$ and suggests studying the manifold $\mathcal T_g^c$ to use Poincar\'e duality. The excerpt omits the exact constructions and maps relating these compact-type spaces, so these statements remain attributed source context rather than newly derived identifications.
\subsection{Short English statement}
Is the anti-invariant part of Torelli-space homology finitely generated after inverting two? The source asks whether the infinite topology is instead accounted for by the Jacobian-locus side and suggests studying compact-type extensions.
\subsection{Sources}
\begin{itemize}
\item Local UnsolvedMath registry, permanent ID 11000088, source code catalogue entry; no proposer or proposal year supplied.
\item Farb, edited volume \emph{Problems on Mapping Class Groups and Related Topics}, chapter context ``Finiteness and Torelli spaces,'' Problem 4.2, PDF page 76; \url{https://www.math.uchicago.edu/~farb/papers/mcgbook.pdf}. The registry records PDF-font-structure extraction and access on 2026-07-29; the volume inventory distinguishes 329 labeled headings by source page.
\item Adjacent local ID 11000087 explicitly supplies the involution eigenspace splitting when two is invertible and the invariant-homology/Jacobian-locus context.
\item The excerpt cites Hamm [9] and [5, p. 152] without full bibliography entries. No exact primary reference for the finite-generation question or its present status is supplied.
\end{itemize}
\end{document}
